\begin{helper}
Let \(V\) be finite, let \(\nu\) be a measure on
\(\operatorname{Finset}(V)\times \mathbb R^V\), and fix an activation fiber
\(A\).  Assume the fixed-fiber cube law and lower-rectangle law used in
\noderef{RandomIndependentSetSampling}.  Let \(\gamma>0\), let
\(a,b,c\in V\) be pairwise distinct, and let \(S_x,S_y,S_z\subseteq V\) be
pairwise disjoint subsets, each disjoint from \(\{a,b,c\}\).  Assume in
addition that the activation fiber has finite mass,
\(\nu[\omega_1=A]<\infty\).  Then the measure of the fiber event with
\(\pi(a)\le\pi(b)\le\pi(c)\), with all \(S_x\)-coordinates below
\(\pi(a)\), all \(S_y\)-coordinates below \(\pi(b)\), and all
\(S_z\)-coordinates below \(\pi(c)\), is at most
\[
\nu[\omega_1=A]\operatorname{ofReal}\left(
\frac{1}{\gamma^3}
\int_{0\le z\le y\le x\le\gamma}
\left(1-\frac{x}{\gamma}\right)^{|S_x|}
\left(1-\frac{y}{\gamma}\right)^{|S_y|}
\left(1-\frac{z}{\gamma}\right)^{|S_z|}\,dx\,dy\,dz\right).
\]
\end{helper}
\begin{proof}
Restrict \(\nu\) to the fiber \(\omega_1=A\).  The cube hypothesis from
\noderef{RandomIndependentSetSampling} says that, inside this fiber, the
activated priority coordinates lie in \([0,1]\) up to a null set.  Since the
fiber mass is finite, the restricted measure on this finite-dimensional cube
is a finite measure.  The lower-rectangle hypothesis gives, for every
\(t_v\in[0,1]\), the value
\[
\nu[\omega_1=A]\operatorname{ofReal}\Bigl(\prod_v t_v\Bigr)
\]
on the box \(0\le \pi(v)\le t_v\) for all \(v\).  Such lower boxes form a
pi-system generating the Borel sigma-algebra of the finite product cube, so
finite-measure uniqueness identifies the restricted priority law with
\(\nu[\omega_1=A]\) times the product of one-dimensional uniform measures.

Under this product law, make the affine change of variables
\[
x=\gamma(1-\pi(a)),\qquad y=\gamma(1-\pi(b)),\qquad
z=\gamma(1-\pi(c)).
\]
Because \(\gamma>0\), the Jacobian contribution is \(\gamma^{-3}\), and
\(\pi(a)\le\pi(b)\le\pi(c)\) becomes \(x\ge y\ge z\) with
\(0\le z\le y\le x\le\gamma\).  The boundary hyperplanes have measure zero,
so they do not affect the value of the integral.

Fix a chamber point \((x,y,z)\).  For \(w\in S_x\) the constraint
\(\pi(w)<\pi(a)=1-x/\gamma\) has one-dimensional uniform measure
\(1-x/\gamma\); the corresponding factors for \(S_y\) and \(S_z\) are
\(1-y/\gamma\) and \(1-z/\gamma\).  The three sets are pairwise disjoint and
are disjoint from \(\{a,b,c\}\), so these outside-coordinate factors multiply
without double-counting and do not interfere with the chamber coordinates.
All remaining coordinates are unrestricted and contribute factor \(1\).
Fubini's theorem for the finite product measure therefore gives exactly the
displayed chamber integral, multiplied by the fiber mass.  Finally, replacing
the original event by its intersection with the full priority cube can only
decrease measure, and the cube complement in the fiber is null, so the claimed
upper bound follows.
\end{proof}
